% ═══════════════════════════════════════════════════════════════════════════
% TAFELBILD: Impuls einführen (DS2)
% DDR-Stil: Kreidetafel-Optik
% ═══════════════════════════════════════════════════════════════════════════

\documentclass[a4paper,landscape,11pt]{article}

\usepackage[utf8]{inputenc}
\usepackage[T1]{fontenc}
\usepackage[ngerman]{babel}
\usepackage[margin=10mm]{geometry}
\usepackage{tikz}
\usepackage{amsmath}
\usepackage{amssymb}

\usetikzlibrary{arrows.meta,decorations.pathmorphing,calc}

% ═══════════════════════════════════════════════════════════════════════════
% FARBEN (Kreide auf Tafel)
% ═══════════════════════════════════════════════════════════════════════════

\definecolor{tafel}{HTML}{1a3d1a}        % Dunkelgrün (DDR-Tafel)
\definecolor{kreide}{HTML}{f5f5dc}       % Beige-Weiß (Kreide)
\definecolor{kreidegelb}{HTML}{ffd700}   % Gelb (Hervorhebung)
\definecolor{kreiderot}{HTML}{ff6b6b}    % Rot (Wichtig)
\definecolor{kreideblau}{HTML}{87ceeb}   % Hellblau (Formeln)

\pagecolor{tafel}
\color{kreide}

\pagestyle{empty}

\begin{document}

\begin{tikzpicture}[remember picture,overlay]
    % Tafelrahmen (Holzoptik)
    \draw[line width=8pt, brown!60!black]
        ([shift={(3mm,3mm)}]current page.south west) rectangle
        ([shift={(-3mm,-3mm)}]current page.north east);
    \draw[line width=3pt, brown!40!black]
        ([shift={(6mm,6mm)}]current page.south west) rectangle
        ([shift={(-6mm,-6mm)}]current page.north east);
\end{tikzpicture}

\vspace*{5mm}

% ═══════════════════════════════════════════════════════════════════════════
% TITEL
% ═══════════════════════════════════════════════════════════════════════════

\begin{center}
{\fontsize{32}{38}\selectfont\bfseries\color{kreidegelb} Der Impuls}\\[3mm]
{\large\color{kreide} Physik Klasse 10 — Bewegungslehre}
\end{center}

\vspace{8mm}

% ═══════════════════════════════════════════════════════════════════════════
% HAUPTINHALT (3 Spalten)
% ═══════════════════════════════════════════════════════════════════════════

\noindent
\begin{minipage}[t]{0.32\textwidth}
\centering

% === LINKE SPALTE: Definition ===
{\Large\bfseries\color{kreidegelb} Was ist Impuls?}\\[5mm]

\begin{tikzpicture}
    % Rahmen
    \draw[kreide, thick, rounded corners=3pt] (0,0) rectangle (6.5,3.5);

    % Text
    \node[text width=6cm, align=center, color=kreide] at (3.25,1.75) {
        {\large Der \textbf{Impuls} beschreibt,}\\[2mm]
        {\large wie schwer es ist,}\\[2mm]
        {\large einen bewegten Körper}\\[2mm]
        {\large zu \textbf{\color{kreiderot}stoppen}.}
    };
\end{tikzpicture}

\vspace{8mm}

{\large\color{kreide} Je größer der Impuls,}\\[2mm]
{\large\color{kreide} desto mehr}\\[2mm]
{\LARGE\bfseries\color{kreiderot} \glqq WUCHT\grqq}\\[2mm]
{\large\color{kreide} hat der Körper.}

\end{minipage}%
\hfill
\begin{minipage}[t]{0.32\textwidth}
\centering

% === MITTLERE SPALTE: Formel ===
{\Large\bfseries\color{kreidegelb} Die Formel}\\[5mm]

\begin{tikzpicture}
    % Doppelter Rahmen (Hervorhebung)
    \draw[kreidegelb, very thick, rounded corners=2pt] (0,0) rectangle (6.5,2.8);
    \draw[kreidegelb, thick, rounded corners=2pt] (0.15,0.15) rectangle (6.35,2.65);

    % Formel
    \node[color=kreideblau] at (3.25,1.4) {
        {\fontsize{36}{42}\selectfont $p = m \cdot v$}
    };
\end{tikzpicture}

\vspace{6mm}

% Legende
\begin{tabular}{@{}r@{\,}c@{\,}l@{}}
{\color{kreideblau}\Large $p$} & {\color{kreide}=} & {\color{kreide}Impuls} \\[3mm]
{\color{kreideblau}\Large $m$} & {\color{kreide}=} & {\color{kreide}Masse} \\[3mm]
{\color{kreideblau}\Large $v$} & {\color{kreide}=} & {\color{kreide}Geschwindigkeit}
\end{tabular}

\vspace{8mm}

% Einheit
{\Large\bfseries\color{kreidegelb} Einheit:}\\[4mm]
{\fontsize{20}{24}\selectfont\color{kreideblau} $[p] = \text{kg} \cdot \dfrac{\text{m}}{\text{s}}$}

\end{minipage}%
\hfill
\begin{minipage}[t]{0.32\textwidth}
\centering

% === RECHTE SPALTE: Beispiele ===
{\Large\bfseries\color{kreidegelb} Beispiele}\\[5mm]

% Tabelle
\begin{tikzpicture}
    % Rahmen
    \draw[kreide, thick] (0,0) rectangle (6.5,5.5);

    % Kopfzeile
    \fill[kreide!20!tafel] (0,4.5) rectangle (6.5,5.5);
    \draw[kreide, thick] (0,4.5) -- (6.5,4.5);

    % Spaltenlinien
    \draw[kreide] (2.2,0) -- (2.2,5.5);
    \draw[kreide] (4.35,0) -- (4.35,5.5);

    % Kopfzeile Text
    \node[color=kreide, font=\small\bfseries] at (1.1,5) {Objekt};
    \node[color=kreide, font=\small\bfseries] at (3.28,5) {m·v};
    \node[color=kreide, font=\small\bfseries] at (5.42,5) {p};

    % Zeilen
    \draw[kreide, thin] (0,3.6) -- (6.5,3.6);
    \draw[kreide, thin] (0,2.4) -- (6.5,2.4);
    \draw[kreide, thin] (0,1.2) -- (6.5,1.2);

    % Daten
    \node[color=kreide, font=\small, align=center] at (1.1,4.05) {Fußball};
    \node[color=kreide, font=\scriptsize] at (3.28,4.05) {0,4·25};
    \node[color=kreiderot, font=\small\bfseries] at (5.42,4.05) {10};

    \node[color=kreide, font=\small, align=center] at (1.1,3) {Sprinter};
    \node[color=kreide, font=\scriptsize] at (3.28,3) {80·10};
    \node[color=kreiderot, font=\small\bfseries] at (5.42,3) {800};

    \node[color=kreide, font=\small, align=center] at (1.1,1.8) {Auto};
    \node[color=kreide, font=\scriptsize] at (3.28,1.8) {1200·30};
    \node[color=kreiderot, font=\small\bfseries] at (5.42,1.8) {36.000};

    \node[color=kreide, font=\small, align=center] at (1.1,0.6) {ICE};
    \node[color=kreide, font=\scriptsize] at (3.28,0.6) {400.000·80};
    \node[color=kreiderot, font=\small\bfseries] at (5.42,0.6) {32 Mio};
\end{tikzpicture}

\vspace{4mm}

{\small\color{kreide} (alle Werte in kg·m/s)}

\end{minipage}

\vspace{8mm}

% ═══════════════════════════════════════════════════════════════════════════
% MERKSATZ (unten)
% ═══════════════════════════════════════════════════════════════════════════

\begin{center}
\begin{tikzpicture}
    % Unterstreichung
    \draw[kreidegelb, very thick, decorate, decoration={snake, amplitude=0.5mm, segment length=3mm}]
        (0,0) -- (24,0);

    % Merksatz
    \node[above=3mm, text width=23cm, align=center, color=kreide] at (12,0) {
        {\Large\bfseries Merke: }
        {\Large Große \textbf{\color{kreiderot}Masse} \textbf{oder} große \textbf{\color{kreiderot}Geschwindigkeit} $\Rightarrow$ großer \textbf{\color{kreidegelb}Impuls}!}
    };
\end{tikzpicture}
\end{center}

\end{document}
